% ============================================================
% Preâmbulo compartilhado — Simulado AF · Microeconomia I (MPE)
% Carregado por simulado-af.tex (\solucaofalse) e
% simulado-af-solucao.tex (\solucaotrue)
% ============================================================
\usepackage[utf8]{inputenc}
\usepackage[T1]{fontenc}
\usepackage[brazil]{babel}
\usepackage{amsmath,amssymb,amsfonts}
\usepackage[a4paper,top=2.2cm,bottom=2.6cm,left=2.2cm,right=2.2cm]{geometry}
\usepackage{graphicx}
\usepackage{fancyhdr}
\usepackage{lastpage}
\usepackage{enumitem}
\usepackage{array}
\usepackage{xcolor}
\usepackage{multirow}

\definecolor{insperred}{HTML}{C8102E}
\definecolor{cinzasol}{gray}{0.93}

\setlength{\parindent}{0pt}
\setlength{\parskip}{0.45em}
\linespread{1.12}

% ---------- header/footer ----------
\pagestyle{fancy}
\fancyhf{}
\lhead{Microeconomia I --- MPE}
\chead{\ifsolucao Simulado 3 AF --- Solução e Rubrica\else Simulado 3 --- Avaliação Final\fi}
\rhead{Página \thepage\ de \pageref{LastPage}}
\renewcommand{\headrulewidth}{0.6pt}
\fancypagestyle{capa}{\fancyhf{}\renewcommand{\headrulewidth}{0pt}}

% ---------- macros de prova ----------
% linha de identificação da capa
\newcommand{\linhaid}[1]{\textbf{#1:}~\hrulefill}

% caixa de resolução (só na versão prova)
\newcommand{\resolucaobox}[1]{%
  \ifsolucao\else
    \par\vspace{0.4em}{\itshape Resolução:}\par\vspace{0.2em}
    \noindent\framebox[\linewidth]{\rule{0pt}{#1}\hfill}\par
  \fi}

% gabarito + solução (só na versão solução) — via comment.sty:
% na versão prova o conteúdo é descartado na leitura (robusto p/ math)
\usepackage{comment}
\ifsolucao
  \specialcomment{solucao}{%
    \par\vspace{0.5em}%
    \begin{list}{}{\leftmargin=0.9em \rightmargin=0.9em}\item[]%
    \textcolor{insperred}{\rule{\linewidth}{0.5pt}}\par\vspace{0.3em}%
  }{%
    \par\vspace{0.2em}\textcolor{insperred}{\rule{\linewidth}{0.5pt}}%
    \end{list}\par\vspace{0.5em}%
  }
\else
  \excludecomment{solucao}
\fi

\newcommand{\gabaritoitem}[1]{\ifsolucao\par\vspace{0.3em}\textbf{Gabarito: #1}\par\fi}

% rótulo de rubrica
\newcommand{\rubricalabel}{\par\vspace{0.35em}\textbf{Rubrica de correção:}\par}

% quadro-resposta de bloco (4 itens objetivos)
\newcommand{\quadroresposta}[1]{%
  \ifsolucao\else
    \par\vspace{0.8em}
    \begin{center}
    \textbf{Quadro-resposta do Bloco #1 --- escreva a alternativa (A--E) ou V/F de cada item:}\\[0.6em]
    \renewcommand{\arraystretch}{1.6}
    \begin{tabular}{|>{\centering\arraybackslash}p{2.6cm}|*{4}{>{\centering\arraybackslash}p{1.4cm}|}}
      \hline
      \textbf{Item} & a) & b) & c) & d) \\ \hline
      \textbf{Resposta} & & & & \\ \hline
    \end{tabular}
    \end{center}
    {\small\itshape Atenção: apenas o quadro acima será considerado na correção dos itens objetivos deste bloco.}\par\vspace{0.6em}
  \fi}
